\section*{Lernziele Selbstkorrigierende Codes}
\begin{todolist}
\item Ich weiss, wie die Prüfziffer von EAN-Codes definiert ist.
\item Ich kann die Prüfziffern eines EAN-Codes berechnen.
\item Ich kann entscheiden, ob ein gegebener EAN-Code korrekt ist.
\item Ich kann bestimmen, welche Fehler (Vertippen einer Ziffer, Vertauschen von zwei Ziffern, Ziffer vergessen) durch eine Prüfziffer erkannt werden.
\item Ich kann den Abstand zweier Bitfolgen gleicher Länge bestimmen.
\item Ich kann den Abstand einer Kodierung bestimmen.
\item Ich kann bestimmen, wie viele Fehler eine Kodierung erkennt.
\item Ich kann bestimmen, wie viele Fehler eine Kodierung korrigieren kann.
\item Ich kann erklären, welchen Abstand \textit{k}-Fehlererkennende Kodierungen haben müssen.
\item Ich kann erklären, welchen Abstand \textit{k}-Fehlererkorrigierende Kodierungen haben müssen.
\item Ich kann den \textit{n}-dimensionalen Hyperwürfel für kleine \textit{n} zeichnen.
\item Ich kann Kodierungen angeben, die einen bestimmten Abstand haben, bzw. eine bestimmte Anzahl Fehler erkennen oder korrigieren können.
\item Ich kann die ``Kartentrick-Kodierung'' einer Bitfolge angeben.
\item Ich kann erklären, wie die optimale Kartentrickkodierung ( = optimale Länge und Breite für das Rechteck) für eine Bitfolge mit einer gegebenen Länge aussieht.
\item Ich kann überprüfen, ob eine gegebene Kartentrick-Kodierung korrekt ist.
\item Ich kann bei einer Kartentrick-Kodierung mit einem Fehler den Fehler finden und korrigieren.
\end{todolist}